\section{L4 SIMD (Flynn / NEON / MISD)}

\subsection{LZ1 Flynn's Taxonomy}

\mult{2}

\begin{definition}{The Four Classes}
    instruction streams $\times$ data streams:
    \begin{itemize}
        \item \textbf{SISD}: 1 thread, 1 pipeline, 1 ALU
        \item \textbf{SIMD}: 1 instr, many FUs/lanes
        \item \textbf{MISD}: lockstep / dependability
        \item \textbf{MIMD}: independent cores
    \end{itemize}
\end{definition}

\begin{concept}{SIMD $\to$ SIMT $\to$ SPMD}
    SIMD is abstracted to \emph{SIMT} (Nvidia, threads in lockstep) and then to the \emph{SPMD} software pattern.
\end{concept}

\multend

\begin{KR}{Recipe: Classify with Flynn}
    Count \emph{instruction} streams $\times$ \emph{data} streams: 1/1 SISD; 1/many SIMD; many/1 MISD; many/many MIMD.
\end{KR}

\begin{example2}{Ex.\ LZ1 Classify these}
    (a) a vector add of 4 lanes; (b) two CPUs running the same code lockstep; (c) a multicore server. Which Flynn class each? What is (a) abstracted to?
     \\ \solution{Solution: }
    (a) \important{SIMD} abstracted to SIMT then SPMD; (b) \important{MISD} (dependability); (c) \important{MIMD}.
\end{example2}

\subsection{LZ2 SPMD Pattern}

\begin{concept}{SPMD}
    \emph{Single Program Multiple Data}: one identical program per UE, each with a unique ID/rank used to branch and index its data. It is the software generalisation of SIMD.
\end{concept}

\begin{KR}{Recipe: Write an SPMD Program}
    Init $\to$ get unique ID/rank $\to$ common code (branch / index by ID) $\to$ distribute data $\to$ collect + combine $\to$ finalise.
\end{KR}

\begin{example2}{Ex.\ LZ2 SPMD vs SIMD}
    Why is SPMD called the software cousin of SIMD?
     \\ \solution{Solution: }
    Both run the \emph{same instructions/program} on \emph{different data}. SIMD does it in hardware lanes (lockstep); SPMD does it across UEs, each picking its data slice by \important{rank} it relaxes the strict lockstep.
\end{example2}

\subsection{LZ3 SIMD Vectorisation Motivation}

\begin{concept}{Why Vectorise}
    DSP kernels FIR filter, Fourier transform, dot product, matrix multiply all reduce to \important{multiply-accumulate (MAC)}. SIMD does one instruction across many lanes in lockstep (e.g.\ a 4-lane parallel add vs a scalar add).
\end{concept}

\begin{KR}{Recipe: Spot a Vectorisable Kernel}
    Is it a \emph{regular, repeated MAC over arrays} with independent elements? $\to$ vectorise. Pointer-chasing / data-dependent control $\to$ does not vectorise well.
\end{KR}

\begin{example2}{Ex.\ LZ3 Which vectorise?}
    FIR filter, dot product, linked-list traversal which suit SIMD and why?
     \\ \solution{Solution: }
    \important{FIR} and \important{dot product} regular MAC over contiguous arrays. \important{Linked-list traversal} does not pointer chasing, no regular data parallelism.
\end{example2}

\subsection{LZ4 NEON Basics (Vector MAC Unit)}

\mult{2}

\begin{definition}{NEON Registers}
    V registers (ARMv8):
    \begin{itemize}
        \item \textbf{Q} = 128-bit; \textbf{D} = 64-bit
        \item lanes: $2{\times}64$ / $4{\times}32$ / $8{\times}16$ / $16{\times}8$
        \item 32-/64-bit IEEE754; mnemonics in the ARM ISA
    \end{itemize}
\end{definition}

\begin{concept}{Vector MAC Unit}
    \begin{itemize}
        \item centres on \important{multiply-accumulate} (VFMA/VFMS)
        \item \important{no divide} only reciprocal estimate (VRSQRTE)
        \item VFP (scalar FP) now deprecated for vectors
    \end{itemize}
\end{concept}

\multend

\begin{definition}{Instruction Format}
    \texttt{V\{mod\}\,op\,\{shape\}\,\{cond\}.dt}:
    \begin{itemize}
        \item \texttt{mod}: Q (saturating), H (halving), D (doubling), R (rounding)
        \item \texttt{shape}: L (long), W (wide), N (narrow) affects result size
        \item \texttt{.dt}: data type (e.g.\ \texttt{.s16})
    \end{itemize}
\end{definition}

\begin{KR}{Recipe: Read a NEON Mnemonic}
    Split \texttt{V} + \texttt{mod} + \texttt{op} + \texttt{shape} + \texttt{.dt}: V $=$ vector op; \texttt{mod}/\texttt{shape} tweak behaviour/width; \texttt{.dt} sets lane type.
\end{KR}

\begin{example2}{Ex.\ LZ4 Decode \texttt{VQADD.s16}}
    What does \texttt{VQADD.s16} do? And why is there no \texttt{VDIV}?
     \\ \solution{Solution: }
    \important{Saturating vector add} of signed 16-bit lanes (Q = saturate instead of wrap on overflow). No \texttt{VDIV}: NEON has \important{no divide} use a reciprocal estimate instead.
\end{example2}

\subsection{LZ5 NEON Intrinsics}

\mult{2}

\begin{definition}{Intrinsics \& Types}
    Compiler-supported C functions (portable vs raw assembly).
    \begin{itemize}
        \item \texttt{uint16x4\_t}: 4$\times$16-bit (64-bit)
        \item \texttt{uint16x8\_t}: 8$\times$16-bit (128-bit)
        \item \texttt{uint16x4x2\_t}: array of two \texttt{uint16x4\_t}
    \end{itemize}
\end{definition}

\begin{concept}{Lanes}
    The individual vector elements. Type + register width define the op: \texttt{vmul\_u16} (4 lanes, D) vs \texttt{vmulq\_u16} (8 lanes, Q).
\end{concept}

\multend

\begin{KR}{Recipe: Decode a NEON Intrinsic}
    Pattern \texttt{v\,op\,q?\,\_\,type}: \texttt{q} present $\Rightarrow$ 128-bit Q (else 64-bit D); \texttt{type} \texttt{<base><bits>x<count>} $\to$ \emph{count} = lanes.
\end{KR}

\begin{example2}{Ex.\ LZ5 (s19) Decode intrinsics}
    What does \texttt{vmulq\_f64} do? How many lanes in \texttt{uint16x4\_t}? What does \texttt{vmul\_lane\_u16(a,b,L)} do?
     \\ \solution{Solution: }
    \texttt{vmulq\_f64}: 128-bit lane-wise multiply of \important{2 doubles} $\to$ \texttt{float64x2\_t}. \texttt{uint16x4\_t}: \important{4 lanes}. \texttt{vmul\_lane}: broadcast \texttt{b[L]} to all lanes, multiply by \texttt{a}.
\end{example2}

\begin{example2}{Ex.\ LZ5 (s20) Lab: vectorise statistics}
    An IoT gateway aggregates a stream (mean, std, skew, kurtosis) to cut comms. Why vectorise, and what is the catch?
     \\ \solution{Solution: }
    These reduce to repeated MAC over the samples $\to$ vectorise with NEON to finish quickly/efficiently. \important{Catch}: they need \emph{divisions} and \emph{floating point}, both expensive on NEON (no divide) compute incrementally and use reciprocal estimates.
\end{example2}

\subsection{LZ6 MISD / Lockstep}

\mult{2}

\begin{concept}{Dependability not Speed}
    MISD spends redundant hardware on \emph{reliability}. Opposite goal to SIMD/MIMD (which buy speed).
\end{concept}

\begin{definition}{Lockstep Variants}
    \begin{itemize}
        \item \textbf{tight}: 2 CPUs 0--2 clk apart, slave compares bus, RESET on diff (ABS)
        \item \textbf{2oo3}: one may fail, re-sync on maintenance
        \item \textbf{loose}: async, \emph{dissimilar} implementations
        \item software pendant: serial or parallel (sync outputs)
    \end{itemize}
\end{definition}

\multend

\begin{KR}{Recipe: Recognise a Dependability Architecture}
    Redundant channels running on the \emph{same} data + compare/vote (RESET on mismatch) $\Rightarrow$ MISD / lockstep. Real example: U96 / Zynq ZU3EG, 2$\times$Cortex-R5F lockstep.
\end{KR}

\begin{example2}{Ex.\ LZ6 ABS controller}
    An ABS unit runs two CPUs that execute the same code and compare every bus access, issuing a RESET on any difference. Flynn class? Purpose?
     \\ \solution{Solution: }
    \important{MISD} (tightly-coupled lockstep). Purpose = \important{dependability}: detect a fault via the mismatch and fail safe (RESET), not to go faster.
\end{example2}

\subsection{Exam FS23 NEON \& Dot Product}

\begin{example2}{Exam FS23 Q2.4 NEON implementation (1P)}
    How is the NEON implemented (RZ/V2L, ARMv8)?
     \\ \solution{Solution: }
    It's a \important{v8 architecture} NEON is integrated into the v8 pipeline (not a separate co-processor).
\end{example2}

\begin{KR}{Recipe: Vectorise a Loop with NEON}
    \paragraph{Load} pack elements into a vector (\texttt{vld1q\_*}).
    \paragraph{Compute} lane-wise op (\texttt{vmulq\_*}, \texttt{vmlaq\_*} for MAC).
    \paragraph{Reduce} horizontal add across lanes (\texttt{vaddvq\_*}) for a dot product.
    \paragraph{Remainder} handle elements beyond a multiple of the lane count with scalar code.
\end{KR}

\begin{example2}{Exam FS23 Q4.4 NEON dot product (6P)}
    Write commented (pseudo/intrinsic) NEON code for the dot product of each row-vector pair (reduced $4\times4$).
     \\ \solution{Solution: }
\begin{lstlisting}[language=C, style=basesmol, aboveskip=-0.05\baselineskip, belowskip=-0.5\baselineskip]
// N=4 fits one int32x4 vector exactly
for (int i = 0; i < N; i++)          // first vector
  for (int j = 0; j < N; j++) {      // second vector
    int32x4_t va = vld1q_s32(&in[i][0]);
    int32x4_t vb = vld1q_s32(&in[j][0]);
    int32x4_t vp = vmulq_s32(va, vb); // lane-wise multiply
    out[i][j] = vaddvq_s32(vp);       // horizontal add of 4 lanes
  }
// edge: if length not a multiple of 4, do the tail scalar
\end{lstlisting}
    Two loops over the pairs; the MAC is one vector op + a horizontal add.
\end{example2}
